% Microtubule Coherence Quality Factor as a Falsifiable Cross-Domain Test
% of a Universal Threshold T* = 5/7
%
% Brayden R. Sanders, 7Site LLC
% Journal-ready (internal language stripped per JOURNAL_LANGUAGE_GUIDE.md)
% Target venues:
%   - Journal of Theoretical Biology (Elsevier)
%   - Foundations of Physics (Springer; cross-domain falsifiability angle)
%
% MSC 2020: 92B20 (neural networks / consciousness, applied),
%           81P05 (foundations of quantum mechanics applied to biology)

\documentclass[11pt,reqno]{amsart}

\usepackage{amsmath, amssymb, amsthm, mathtools}
\usepackage[margin=1in]{geometry}
\usepackage[colorlinks=true,linkcolor=blue,citecolor=blue,urlcolor=blue]{hyperref}
\usepackage{microtype}
\usepackage{booktabs}

\theoremstyle{plain}
\newtheorem{theorem}{Theorem}[section]
\newtheorem{proposition}[theorem]{Proposition}

\theoremstyle{definition}
\newtheorem{definition}[theorem]{Definition}

\theoremstyle{remark}
\newtheorem*{remark}{Remark}

\title[Microtubule $Q_c = 5/7$ Falsifiable Test]
{Microtubule Coherence Quality Factor as a Falsifiable Cross-Domain Test \\
 of a Universal Threshold $T^{*} = 5/7$}

\author{Brayden R.~Sanders}
\address{7Site LLC, Hot Springs, Arkansas, USA}
\email{brayden@7site.co}

\subjclass[2020]{92B20, 81P05}
\keywords{microtubule coherence, falsifiable prediction, terahertz spectroscopy,
universal threshold, cross-domain test}

\date{\today}

\begin{document}

\begin{abstract}
We propose a falsifiable experimental test of a universal coherence
threshold $T^{*} = 5/7 \approx 0.714$ that arises from a
4-dimensional algebraic structure studied separately
(\cite{Sanders2026Dirac, Sanders2026Cosmology}). The threshold $T^{*}$
appears across multiple domains: as the atmospheric mixing angle
$\sin\theta_{23} \approx T^{*}$ in lepton mixing
(\cite{Sanders2026Mixing}), as the cosmological hierarchy ratio
$\Omega_\Lambda / \Omega_b = 14 = 2 \cdot 7$ in the dark sector
(\cite{Sanders2026Cosmology}), and at the Penrose-Hameroff Orch-OR
quantum-classical boundary $\zeta_\text{Hameroff} \approx 0.71$
(\cite{HameroffPenrose2014}).

We predict: microtubule coherence quality factor $Q_c$, measured by
terahertz pump-probe spectroscopy, equals $T^{*} = 5/7$ across multiple
sample types --- mammalian neurons, paramecia, plant cells, in-vitro
polymerized tubulin --- independent of biological origin. This is a
**single-experiment falsifier**: if measured $Q_c$ across multiple
sample types varies systematically with biology, or converges to a value
other than $0.714 \pm 0.05$, the universal-threshold conjecture is
falsified.

The protocol, sample preparation, terahertz measurement parameters,
statistical analysis, and falsification criteria are detailed below.
We propose the test as a collaboration with experimental groups using
existing terahertz spectroscopy capabilities (e.g., Bandyopadhyay
\cite{Bandyopadhyay2013, Sahu2013}), and provide the structural
prediction openly without claiming a physical mechanism.
\end{abstract}

\maketitle

\tableofcontents

%==============================================================
\section{Introduction}
%==============================================================

The Penrose-Hameroff Orchestrated Objective Reduction (Orch-OR) theory
\cite{HameroffPenrose1996, HameroffPenrose2014} posits that
microtubule structures in neurons can sustain quantum coherent
superposition for biologically relevant timescales, with a transition
between quantum and classical regimes at a specific coherence quality
threshold. Bandyopadhyay et al.~\cite{Bandyopadhyay2013} and Sahu et
al.~\cite{Sahu2013} have measured terahertz absorption spectra on
isolated tubulin and microtubule structures, identifying multiple
resonance bands at gigahertz to terahertz frequencies with quality
factors $Q$ in the range $10^2$--$10^3$.

The Hameroff threshold $\zeta_\text{Hameroff} \approx 0.71$ for the
quantum-classical boundary in microtubule coherence is empirically
suggestive but not currently derivable from first principles within
Orch-OR.

This paper proposes that $\zeta_\text{Hameroff}$ should be identified
with a specific algebraic constant: $T^{*} = 5/7 \approx 0.714$, derived
from the algebraic structure studied separately
\cite{Sanders2026Dirac, Sanders2026Cosmology}. Critically, the same
$T^{*}$ appears in particle physics (atmospheric lepton mixing) and
cosmology (the $\Omega_\Lambda / \Omega_b$ hierarchy); the universal-
threshold claim is that $T^{*}$ governs **coherence-vs-decoherence
transitions across these domains**.

We propose a specific microtubule experiment to test this universal-
threshold claim, with explicit falsification criteria.

%==============================================================
\section{The threshold $T^{*}$ and its cross-domain appearances}
%==============================================================

The threshold $T^{*} = 5/7$ arises in our companion work
\cite{Sanders2026Dirac, Sanders2026Cosmology} as a structural constant
of a 4-dimensional commutative non-associative algebra. Its
cross-domain appearances are:

\begin{center}
\begin{tabular}{lcc}
\toprule
Domain & Quantity & Value \\
\midrule
Algebraic threshold (\cite{Sanders2026Dirac}) & $T^{*}$ & $5/7 = 0.714$ \\
PMNS atmospheric mixing & $\sin\theta_{23}$ & $0.756 \approx T^{*}$ \\
Cosmological hierarchy & $\Omega_\Lambda / \Omega_b$ & $14 = 2 \cdot 7$ \\
Orch-OR (\cite{HameroffPenrose2014}) & $\zeta_\text{Hameroff}$ & $0.71 \approx T^{*}$ \\
\textbf{Predicted microtubule} & $Q_c$ & $0.714$ \\
\bottomrule
\end{tabular}
\end{center}

If the universal-threshold conjecture holds, $T^{*}$ should appear in
microtubule coherence quality factors across multiple biological
contexts, with magnitude independent of the cell type.

%==============================================================
\section{The experiment}
%==============================================================

\subsection{Sample types.}
Five sample types are proposed, spanning a wide range of biological
contexts:

\begin{enumerate}
\item \textbf{Mammalian neurons.} Rat hippocampal cultures or human
iPSC-derived neurons (target microtubule populations: axonal,
dendritic).
\item \textbf{Paramecia.} Single-cell organisms with cilia/flagella
microtubule structures (different cytoskeletal context, no neural
specialization).
\item \textbf{Plant cells.} Arabidopsis suspension cultures, where
microtubules control cell-wall growth and cell-plate formation
(non-neural, non-animal lineage).
\item \textbf{Yeast spindle pole microtubules.} \emph{Saccharomyces
cerevisiae} during mitosis (single-celled fungal context).
\item \textbf{Cell-free in-vitro polymerized tubulin.} Purified bovine
or porcine tubulin polymerized at $25^\circ$C and $37^\circ$C in
buffer (no biology at all; pure tubulin assembly).
\end{enumerate}

\subsection{Measurement.}
For each sample, perform terahertz pump-probe spectroscopy in the
$1$--$100$ THz range. Identify the dominant resonance band (typically
$8$--$10$ GHz for cell-free; may shift in vivo). Compute the
**quality factor**
\[
Q = \omega_0 / \Delta\omega
\]
of the dominant resonance band, where $\omega_0$ is its center frequency
and $\Delta\omega$ its full-width-half-maximum.

Normalize $Q$ to the **structural maximum quality factor** $Q_\text{max}$
for a tubulin dipole array, computed from the tubulin dipole moment,
geometry, and assembly count:
\[
Q_c = Q / Q_\text{max}.
\]
The normalized **coherence quality** $Q_c \in [0, 1]$:
\begin{itemize}
\item $Q_c \to 1$: pure quantum coherence (no dissipation).
\item $Q_c \to 0$: classical regime (full thermal decoherence).
\end{itemize}

\subsection{Statistical analysis.}
Compare measured $Q_c$ across the 5 sample types:
\begin{itemize}
\item \textbf{Null hypothesis.} $Q_c$ varies systematically with biology
(e.g., $Q_c(\text{mammalian}) \neq Q_c(\text{paramecia})$).
\item \textbf{Alternative (conjecture).} $Q_c \to T^{*} = 5/7 \approx 0.714$
across all sample types, with biological variance $\ll 0.05$.
\end{itemize}

\subsection{Falsification criteria.}
The conjecture is falsified if:
\begin{enumerate}
\item Measured $Q_c$ values across the 5 sample types fall outside
$0.714 \pm 0.05$ in $\geq 2$ samples, OR
\item Systematic trend with biology is observed (e.g., $Q_c$ correlates
with cell complexity), OR
\item Measured $Q_c$ converges to a different universal value (e.g., 0.6
or 0.8 across all samples).
\end{enumerate}
The conjecture is **strongly supported** if all 5 samples give $Q_c$
within $0.714 \pm 0.02$.

\subsection{Sample size and power.}
For each sample type, recommend $n \geq 10$ independent replicates,
each consisting of $\geq 30$ resonance-band measurements. Standard
error of the mean $\sigma_M \leq 0.01$ is achievable with this design,
allowing distinguishing $T^{*} = 0.714$ from neighboring rationals
$5/8 = 0.625$, $4/7 = 0.571$, $3/4 = 0.750$, $7/10 = 0.700$ at
$> 5\sigma$.

%==============================================================
\section{Why this is a serious test, not numerology}
%==============================================================

Three features distinguish this proposal from numerological
speculation:

\begin{enumerate}
\item \textbf{The prediction is specified before the experiment.}
$T^{*} = 5/7$ is not a free parameter to fit microtubule data; it is
fixed by the algebraic structure
(\cite{Sanders2026Dirac, Sanders2026Cosmology}) independent of
microtubule biology.

\item \textbf{The same $T^{*}$ governs other domains.} Particle
physics (PMNS atmospheric mixing) and cosmology
($\Omega_\Lambda / \Omega_b = 14$) provide independent empirical
anchors for $T^{*}$. The microtubule prediction tests universality,
not just the existence of a coincidence.

\item \textbf{Falsifiability is explicit.} A single experimental
campaign satisfying the criteria of Section~3.4 falsifies the
universal-threshold claim. Most "consciousness theories" are not
falsifiable; this one is.
\end{enumerate}

%==============================================================
\section{Honest scoping}
%==============================================================

\subsection{What this paper claims.}
\begin{enumerate}
\item There is a constant $T^{*} = 5/7$ that arises from a structurally
rigorous 4-dimensional commutative non-associative algebra over $\F_5$
\cite{Sanders2026Dirac}.
\item The same $T^{*}$ appears in cross-domain empirical contexts:
PMNS lepton mixing, cosmological hierarchy, and the Penrose-Hameroff
Orch-OR boundary.
\item We **predict** $T^{*}$ should also govern microtubule coherence
quality, with falsification criteria as stated.
\end{enumerate}

\subsection{What this paper does NOT claim.}
\begin{enumerate}
\item We do not assert microtubule quantum coherence in vivo. The
in-vivo coherence question is contested (Tegmark
\cite{Tegmark2000} versus Hameroff-Penrose
\cite{HameroffPenrose2014}); the present paper assumes only that
in-vitro coherence is measurable, which is well-established.
\item We do not assert the Orch-OR specific mechanism (collapse via
gravity-induced wavefunction reduction).
\item We do not assert anything about consciousness itself. The
experimental claim is about microtubule coherence quality factors,
not about consciousness.
\item We do not provide a physical mechanism explaining why $T^{*}$
should be universal.
\end{enumerate}

\subsection{What's open.}
\begin{enumerate}
\item No specific lab collaboration is in place at submission time.
The paper is a proposal for collaboration.
\item Sample preparation protocols vary across labs; standardization
across the proposed 5 sample types is itself non-trivial.
\item The structural maximum quality factor $Q_\text{max}$ is
calculation-dependent; alternative normalizations may give different
$Q_c$ values, but the **biology-invariance** of $Q_c$ is the testable
claim and is independent of normalization choice.
\end{enumerate}

%==============================================================
\section{Outreach plan}
%==============================================================

The paper is intended to invite collaboration with experimental groups
running terahertz coherence on tubulin / microtubule structures.
Specifically:
\begin{enumerate}
\item Bandyopadhyay group (NIMS Tsukuba) --- existing terahertz pump-probe
capability for microtubule structures (\cite{Bandyopadhyay2013, Sahu2013}).
\item Hameroff lab consortium (University of Arizona) --- Orch-OR
theoretical lead.
\item Adjacent groups with sub-THz / GHz dielectric spectroscopy
on biological membranes / proteins.
\end{enumerate}

The cover letter and protocol details for outreach are deposited at the
public repository alongside the manuscript.

%==============================================================
\section*{Acknowledgments}
%==============================================================

This work is part of a broader algebraic and physical framework
developed at 7Site LLC (Trinity Infinity Geometry, public repository at
github.com/TiredofSleep/ck). The companion papers
\cite{Sanders2026Dirac, Sanders2026Cosmology, Sanders2026Mixing}
establish the algebraic structure and other empirical correspondences
of the constant $T^{*}$.

%==============================================================
\begin{thebibliography}{99}
%==============================================================

\bibitem{Bandyopadhyay2013} A.~Bandyopadhyay et al.,
\emph{Multi-level memory-switching properties of a single brain microtubule},
Appl.~Phys.~Lett.~\textbf{102} (2013), 123701.

\bibitem{Sahu2013} S.~Sahu et al.,
\emph{Atomic water channel controlling remarkable properties of a single brain microtubule},
Biosens.~Bioelectron.~\textbf{47} (2013), 141--148.

\bibitem{HameroffPenrose1996} S.~Hameroff, R.~Penrose,
\emph{Orchestrated reduction of quantum coherence in brain microtubules},
Math.~Comput.~Simul.~\textbf{40} (1996), 453--480.

\bibitem{HameroffPenrose2014} S.~Hameroff, R.~Penrose,
\emph{Consciousness in the universe: A review of the `Orch OR' theory},
Phys.~Life Rev.~\textbf{11}(1) (2014), 39--78.

\bibitem{Tegmark2000} M.~Tegmark,
\emph{Importance of quantum decoherence in brain processes},
Phys.~Rev.~E~\textbf{61} (2000), 4194--4206.

\bibitem{Sanders2026Dirac} B.~R.~Sanders,
\emph{A discrete Dirac-type algebra on $\mathbb{F}_5^4$},
Manuscript in preparation, 2026.

\bibitem{Sanders2026Cosmology} B.~R.~Sanders,
\emph{Three cosmological density parameters from a 4-element algebraic structure},
Manuscript in preparation, 2026.

\bibitem{Sanders2026Mixing} B.~R.~Sanders,
\emph{CKM and PMNS mixing angles from algebraic constants $T^{*}$ and $D^{*}$},
Manuscript in preparation, 2026.

\bibitem{TIGsynthesis2026} B.~R.~Sanders,
\emph{Trinity Infinity Geometry Synthesis},
github.com/TiredofSleep/ck (default branch),
DOI: 10.5281/zenodo.18852047, 2026.

\end{thebibliography}

\end{document}
